%% =====================================================================
\chapter{Introduction}\label{ch:intro}

\section{What MediChain is}

MediChain is a web platform that runs the day-to-day information flow of a
hospital for four kinds of people:

\begin{description}[leftmargin=2.4cm,style=nextline]
  \item[Patients] book appointments, read their own medical records, decide which
    doctors may see those records and for how long, follow their medications and
    home-monitoring vital signs, pay bills and file insurance claims.
  \item[Doctors] manage their schedule, ask patients for consent, write and amend
    medical records, attach files such as lab reports, prescribe, watch live vital
    signs from patients' devices, and in a genuine emergency open a record without
    consent (``break-glass''), which is logged and reported.
  \item[Administrators] verify the credentials of new doctors and insurance agents,
    manage user accounts and IoT devices, configure the hospital, read the security
    audit trail and check the integrity of the whole database.
  \item[Insurance agents] work through a queue of claims, approve or deny them,
    record settlements and read financial reports.
\end{description}

Two ideas run through the whole system and make it more than a CRUD application:

\begin{enumerate}
  \item \textbf{Zero trust.} No request is trusted because of where it comes from
    or because the user logged in earlier. Every API call is authenticated, scored,
    checked against the user's role, approval status and (for medical data) the
    patient's consent, and the decision is recorded.
  \item \textbf{Tamper evidence.} Every change to clinical or financial data, and
    every access decision, is written to a hash chain: a sequence of entries where
    each one contains the SHA-256 hash of the previous one. If anyone edits the
    database directly, bypassing the application, the system can prove it.
\end{enumerate}

\section{The problem it addresses}

Hospital data has three properties that make ordinary web-application practice
insufficient:

\begin{itemize}
  \item \textbf{It is extremely sensitive.} A diagnosis or a prescription must be
    visible only to the patient and the clinicians the patient has chosen.
    ``Logged-in user with the doctor role'' is not a good enough rule.
  \item \textbf{It must be trustworthy.} Insurance decisions, legal disputes and
    clinical care all depend on records being exactly what the doctor wrote. A
    silently edited record is worse than a missing one.
  \item \textbf{Emergencies happen.} An unconscious patient cannot grant consent.
    The system must allow emergency access without making it invisible.
\end{itemize}

The original MediChain research paper proposed Hyperledger Fabric (a blockchain)
and IPFS (distributed file storage) to meet these needs. This implementation
keeps the guarantees and replaces the heavy infrastructure with a hash-chain
ledger and file hashes stored in MongoDB (Chapter~\ref{ch:architecture} explains
the reasoning).

\section{What changed from the starting point}

When this work began, the repository contained a polished but static frontend
(every number on every dashboard was hardcoded in JavaScript) and one backend
service that handled login only. Table~\ref{tab:before-after} summarises what
exists now.

\begin{table}[H]
\centering\small
\caption{Starting point versus completed system}\label{tab:before-after}
\begin{tabularx}{\textwidth}{@{}P{3.4cm}YY@{}}
\toprule
\textbf{Area} & \textbf{Before} & \textbf{Now} \\
\midrule
Backend & One auth service, 6 endpoints, a broken logout & One API service, 68 endpoints in 13 route files \\
Frontend data & Hardcoded arrays, simulated login & Every screen loads from the API; real loading, empty and error states \\
Login & Simulated in the browser & Password + emailed OTP, JWT with rotating refresh cookie \\
Appointments, records, claims, admin & Not implemented & Implemented, tested, wired to all four dashboards \\
Ledger & A JSON-lines audit file & Two hash chains in MongoDB, full verification, tamper demo \\
Files & None & GridFS storage, SHA-256 anchored in the ledger, consent-gated downloads \\
IoT vitals & Static numbers & Device API keys, ingestion endpoint, live monitor, alerts, simulator \\
Tests & None & 105 backend + 20 frontend tests, about 93\% backend line coverage \\
Deployment & None & Dockerfile, docker-compose, GitHub Actions CI \\
\bottomrule
\end{tabularx}
\end{table}

\section{How to read this document}

\begin{itemize}
  \item Chapters~\ref{ch:requirements}--\ref{ch:setup} are the overview:
    requirements, architecture and how to run the project. Read these first.
  \item Chapters~\ref{ch:foundation}--\ref{ch:admin} walk through the backend one
    module at a time. Each section states \emph{what} the module is for,
    \emph{why} it is built that way, and then shows the real code with an
    explanation of each part.
  \item Chapter~\ref{ch:frontend} does the same for the browser code.
  \item Chapters~\ref{ch:reference}--\ref{ch:demo} are reference material: data
    model, full API list, tests, deployment, the demo script, the security model
    and what is left to do.
\end{itemize}

Conventions: file paths are relative to the repository root unless they start
with \file{src\codesl{}}, which means \file{backend\codesl{}services\codesl{}api\codesl{}src}.
Code listings are exact copies; where a listing omits lines, the gap is marked
with a \code{\codesl{}\codesl{} ...} comment.
